\documentclass[12pt,a4paper]{article}

\usepackage[spanish,es-nodecimaldot]{babel}
\usepackage[utf8]{inputenc}
\usepackage[T1]{fontenc}
\usepackage{lmodern}
\usepackage{amsmath,amssymb}
\usepackage{graphicx}
\usepackage{booktabs}
\usepackage{hyperref}
\usepackage{enumitem}
\usepackage{geometry}
\usepackage{listings}
\usepackage{xcolor}

\geometry{margin=2.5cm}

\lstset{
  basicstyle=\footnotesize\ttfamily,
  breaklines=true,
  frame=single,
}

\begin{document}
\section{01 --- Fundamentos y objetivo del TFG}

\subsection{1) Qué problema resuelve el proyecto}

El sistema toma decisiones binarias sobre flujos de red:

\begin{itemize}
  \item \texttt{0 = PERMIT}
  \item \texttt{1 = BLOCK}
\end{itemize}

La idea central es que en ciberseguridad los errores no cuestan lo mismo:

\begin{itemize}
  \item \textbf{Falso negativo (FN)}: dejar pasar un ataque suele ser el error más grave.
  \item \textbf{Falso positivo (FP)}: bloquear tráfico benigno también duele, pero normalmente menos.
\end{itemize}

Por eso el proyecto se formula con \textbf{Aprendizaje por Refuerzo (RL)} en lugar de quedarse solo
en clasificación supervisada clásica.

\subsection{2) Objetivo real del trabajo}

El objetivo no es solo ``sacar accuracy alta'', sino construir un pipeline reproducible de extremo a
extremo:

\begin{enumerate}
  \item adaptar datos de red a un contrato común
  \item entrenar un agente RL sobre ese contrato
  \item validar con pruebas anti-autoengaño (A/B/C + leave-one-exact-CSV-out)
  \item llevar el modelo a inferencia offline en tráfico real de laboratorio
\end{enumerate}

\subsection{3) Estructura por fases}

\subsubsection{Fase 1 (más madura)}

\begin{itemize}
  \item entrenamiento y validación offline
  \item dataset principal: CICIDS2017
  \item baseline histórico: NSL-KDD
\end{itemize}

\subsubsection{Fase 2 (implementada pero abierta en robustez)}

\begin{itemize}
  \item inferencia offline sobre flujos extraídos de tráfico de laboratorio
  \item script mantenido: \verb|scripts/predict_real_traffic_v2.py|
  \item foco actual: controlar el \textbf{domain shift} entre dataset y tráfico real
\end{itemize}

\subsection{4) Qué está implementado y qué no}

\subsubsection{Implementado}

\begin{itemize}
  \item esquema canónico fijo de 76 features
  \item máscara de missingness (otras 76)
  \item observación final de 152 dimensiones
  \item entorno RL custom con Gymnasium
  \item entrenamiento QRDQN
  \item validaciones A/B/C
  \item validación leave-one-exact-CSV-out en código
  \item inferencia robusta offline en Phase~2
\end{itemize}

\subsubsection{No implementado todavía}

\begin{itemize}
  \item bloqueo activo en tiempo real (\verb|iptables|/\verb|nftables|)
  \item cierre completo de calibración en tráfico benigno real
  \item despliegue productivo
  \item línea multiagente/adversarial
\end{itemize}

\subsection{5) Mensaje fuerte para defensa}

El valor del TFG está en combinar:

\begin{itemize}
  \item rigor de ingeniería de datos (contrato de observación estable)
  \item diseño de decisión con costes asimétricos (reward shaping)
  \item validación que evita sobreinterpretar métricas cómodas
  \item transición realista hacia tráfico real, con límites reconocidos
\end{itemize}

\end{document}
